\begin{examplebox}
\textbf{\textcolor{CExample}{例 1（基础）}}

\textbf{\textcolor{CExample}{题目：}} 在非直角三角形 $\triangle ABC$ 中，已知 $\tan A = 2$，$\tan B = 3$，求 $\tan C$ 的值。

\textbf{\textcolor{CExample}{解题步骤：}}
\begin{description}[style=nextline, leftmargin=2.8em, labelsep=0.5em, itemsep=0.45em]
    \item[\textbf{第一步（恒等式列写）：}] 由结论得
    \[
    \tan A + \tan B + \tan C = \tan A \cdot \tan B \cdot \tan C.
    \]
    \par\textbf{理由：}
    三角形非直角，恒等式成立。

    \item[\textbf{第二步（反解 $\tan C$）：}] 视 $\tan C$ 为未知数，移项整理
    \[
    \tan C(1 - \tan A \tan B) = -(\tan A + \tan B),
    \]
    即
    \[
    \tan C = \frac{\tan A + \tan B}{\tan A \tan B - 1}.
    \]
    \par\textbf{理由：}
    将乘积项移至一边，提取公因子 $\tan C$，分母不为零。

    \item[\textbf{第三步（代入求值）：}] 代入 $\tan A = 2$，$\tan B = 3$：
    \[
    \tan C = \frac{2+3}{2\times 3 - 1} = \frac{5}{5} = 1.
    \]
    \par\textbf{理由：}
    直接计算，结果 $\tan C = 1$ 对应角 $C = \dfrac{\pi}{4}$。
\end{description}

\textbf{\textcolor{CExample}{关键结论：}} \textbf{$\tan C = 1$，且验证 $\tan A+\tan B+\tan C = 6 = 2\times 3\times 1$，与恒等式一致。}

\medskip
\textbf{\textcolor{CExample}{例 2（稍变形）}}

\textbf{\textcolor{CExample}{题目：}} 已知非直角三角形 $\triangle ABC$ 满足 $\tan A + \tan B + \tan C = \tan A\tan B\tan C$，且三者之和为 $-3$，试判断三角形的形状。

\textbf{\textcolor{CExample}{解题步骤：}}
\begin{description}[style=nextline, leftmargin=2.8em, labelsep=0.5em, itemsep=0.45em]
    \item[\textbf{第一步（等量替换）：}] 由恒等式及条件得
    \[
    \tan A\tan B\tan C = \tan A+\tan B+\tan C = -3.
    \]
    \par\textbf{理由：}
    恒等式保证和与积相等，可直接代换。

    \item[\textbf{第二步（符号分析）：}] 乘积 $-3 < 0$，说明三个正切中必有一个为负值。
    \par\textbf{理由：}
    三角形内角的正切：钝角的正切为负，锐角的正切为正；乘积为负当且仅当有奇数个负因子。

    \item[\textbf{第三步（形状判定）：}] 三个角中恰有一个钝角（不可能有两个钝角），故 $\triangle ABC$ 为钝角三角形。
    \par\textbf{理由：}
    三角形内角和 $\pi$ 至多允许一个钝角，负乘积直接推出存在钝角。
\end{description}

\textbf{\textcolor{CExample}{关键结论：}} \textbf{三角形为钝角三角形，无需逐角求解，仅用符号即可判定。}
\end{examplebox}
